% !TeX root = ../papers/nc-p3.tex
\section{二重点の六次条件が独立な辺の正則延長}
\label{sec:opposite-sextic-threshold}
\begin{theorem}\label{thm:opposite-sextic-threshold}
四辺有限な標準有限Frobenius窓で、一位置の長さ二十の射影を $X$ とする。
反対二位置が添字の標識を保って射影同値でなければ
\[
                  h^0(\mathcal I_X^2(6))\ge5 .
\]
特に二重化の長さ80の条件が六次で独立なら、反対交換は対称である。
$X$ の被約性は必要ない。
\end{theorem}
\begin{proof}
反対位置の二乗を最初の二次式で読んだ対称行列を $N(x)$ とする。
その十成分は全二次式の基底であり、$X$ 上では単元倍を除き反対座標の平方である。
単元の対角成分を局所消去すると三次小行列式は全て $\mathcal I_X^2$ に入る。
その線型包が四次元以下と仮定する。一般階数1は同じVeroneseを与え、
2は定理\ref{thm:nonreduced-opposite-root-separation}で排除済みである。

一般階数3なら $\operatorname{adj}N=hgg^{\mathsf t}$ と原始に書き、
$g$ の次数を $e$、その成分の張る次元を $r$ とする。
$e=0$ または $r=1$ なら核が固定される。積 $g_i g_j$ の張る次元は $2r-1$ 以上である。
単項式順序で相異なる先頭項を持つ基底を取り、最小のものと全基底の積、
および最大のものと最小以外の積を取ればよい。従って $r=2$ である。
合同変換で $g=(a,b,0,0)^{\mathsf t}$ とし、$a,b$ は互いに素である。
各行の先頭二成分は $(b,-a)$ の $2-e$ 次倍となる。
$e>2$ は不可能で、$e=2$ でも先頭二列の七成分が二次元以下に入り矛盾する。
$e=1$ なら $Ng=0$ を $X$ 上で読むと最初と反対の文字の非零な双一次関係である。
定理\ref{thm:finite-opposite-product-gap}により反対二位置が同定される。

一般階数4とする。逆行列を射影的に取る写像は稠密開で対合であり、
$x\mapsto[\operatorname{adj}N(x)]$ は像と双有理である。
従って像の閉包は四次元線型包の射影空間全体で、その線型表示を $L(y)$ とする。
十成分の基底変換 $T$ により $N(x)=T(xx^{\mathsf t})$ と書く。
$T^{-1}(\operatorname{adj}L(y))$ は恒等的に階数1の三次式行列なので、
一意分解で $h(y)u(y)u(y)^{\mathsf t}$ となる。
$[u(y)]$ は双有理だから成分は四つの独立な一次式、$h$ は非零一次式である。
$u$ を座標として $\operatorname{adj}L(x)=h(x)N(x)$ を得る。

平面 $h=0$ の全点で $L$ の階数はちょうど2である。実際1以下なら
三次小行列式の一次微分まで全て零であり、$h$ に横断的な微分で $N$ の基点が生じる。
この平面の階数2二次形式は二つの相異なる一次式の積である。
二因子の順序付けは二次エタール被覆をなし、射影平面上では分裂する
（$2$ は可逆でPicard群に二捩れがなく、定数体は代数閉）。
従って二つの因子はそれぞれ射影空間への射を与える。
その引戻し次数は非負で和が1だから、一方は固定で他方は一次である。
動く因子の三次元空間は固定因子を含まない。含めば平方が現れるからである。
従って合同変換と一次の座標変更の後
\[
            L=\begin{pmatrix} ch & v^{\mathsf t}\\v & hE \end{pmatrix},
\]
ここで $(h,v_1,v_2,v_3)$ は座標、$c,E$ は定数である。
$N$ の左上成分が $h^2\det E$ なので $E$ は可逆であり、合同で $E=I_3$ に取れる。
余因子を展開すると
\[
 N=ss^{\mathsf t}+q C,\quad s=(h,-v)^{\mathsf t},\quad
 q=ch^2-v^{\mathsf t}v,\quad C=\operatorname{diag}(0,1,1,1).
\]
$q\ne0$ の各点では階数が2以上である。$q=0$ 上でも、全二次小行列式は $q$ 倍であり、
各点でその商の一つは単元である。実際 $C$ が $ss^{\mathsf t}$ での階数1の接空間に入れば、
三次元空間 $s^\perp$ が $C$ に対して全等方となるが、階数3の四変数形式にはそのような空間はない。
よって $X$ 上での階数1からスキームとして $q|_X=0$ となる。
二次評価の単射性に反してこの場合も起こらない。
\end{proof}

\begin{corollary}\label{cor:sextic-threshold-finite-extension}
指定した有限辺の第一射影が $h^0(\mathcal I_X^2(6))=4$ を満たすとき、
標準正則延長は添字保持同型まで高々一つで、その存在は有限判定できる。
特に六個の双一次関係の一般の六次元空間は標準正則延長を持たない。
\end{corollary}
\begin{proof}
まず任意の標準正則延長の四辺は全て有限である。実際どこかの射影像に面があれば、
定理\ref{thm:no-moving-plane-pencils}と定理\ref{thm:isolated-surface-fibers}の
面の列は全位置で続き、指定辺の有限性に反する。全射影像が高々一次元なら
定理\ref{thm:unrestricted-corridor-exclusion}により分岐はなく、
命題\ref{prop:surface-degree-bound}により全辺は同じスキームのグラフである。
従って上の定理により次の辺は指定辺の転置である。系\ref{cor:octic-free-finite-extension}の
三次交叉からの復元と一交換対称窓の正則判定は、この同定後は被約性を使わない。
一般辺の第一射影は有限被約である。再現計算
\path{computations/verify_opposite_sextic_threshold.py} では一つの整数係数辺を使い、
射影の最大四次小行列式で割った五次商における六次形式の四偏導の評価階数80、
命題\ref{prop:finite-flip-obstructions}の交代化の階数24を良い素数で確認する。
一般辺では最大小行列式が期待余次元で五次商は二十次元であるため、
これらは標数零の一般点での最大階数を与える。最初の階数から六次核は4、
後者から転置辺の三次交叉は零となり、必要な四次元に達しない。
\end{proof}
